\section{The Hodge structure generated by the top forms}\label{sec:hodge}

We will bound the order of a birational map on a one-dimensional space of
pluricanonical forms by realizing its scalar as an eigenvalue of an integral
Hodge structure of bounded rank.  The relevant Hodge structure is usually
much smaller than the full middle cohomology.  We first recall its basic
properties and prove the metric estimate that will make its rank accessible
in Section~\ref{sec:diagonal}.

For a smooth projective complex $n$-fold $U$, let
\[
 T(U)\subset H^n(U,\Q)
\]
be the smallest rational Hodge substructure whose complexification contains
$H^{n,0}(U)$, and put $r(U)=\dim_\Q T(U)$.  Such a smallest substructure
exists: intersections are Hodge substructures, and finite dimensionality
reduces an arbitrary intersection to a finite one.

\begin{lemma}\label{hodge:basic}
The Hodge structure $T(U)$ is birationally invariant.  Cup product with
any rational divisor class annihilates $T(U)$; in particular, $T(U)$ is
primitive for every rational ample class.  Moreover, if $f:Z\to U$ is a
morphism from a smooth projective variety with $\dim f(Z)<n$, then
$f^*T(U)=0$.
\end{lemma}

\begin{proof}
For a birational morphism $p:W\to U$ between smooth projective varieties,
$p^*:H^n(U,\Q)\to H^n(W,\Q)$ is an injective Hodge morphism, with left
inverse $p_*$.  Birational invariance of holomorphic top forms gives
$H^{n,0}(W)=p^*H^{n,0}(U)$.  The image $p^*T(U)$ therefore contains
$T(W)$ by minimality.  Its inverse image of $T(W)$ is a Hodge substructure
containing $H^{n,0}(U)$, and hence equals $T(U)$.  Thus
$T(W)=p^*T(U)$.  Common resolutions give the assertion for birational maps.

The kernel of cup product with a rational divisor class is a rational
Hodge substructure, with the usual Tate twist in the target.  It contains
$H^{n,0}(U)$ because $H^{n+1,1}(U)=0$, so it contains $T(U)$.  Finally,
$f^*$ vanishes on holomorphic top forms when its image has dimension less
than $n$.  Its kernel is again a rational Hodge substructure, which proves
the last assertion.
\end{proof}

For a normal projective variety $V$, we also write $r(V)=r(U)$, where
$U$ is any smooth projective resolution. Lemma~\ref{hodge:basic} makes
this independent of the choice of $U$.

We use the Hodge norm $\norm{z}_\omega$ of a cohomology class on a compact
K\"ahler manifold: this is the $L^2$-norm of its $\omega$-harmonic
representative.  It is also the minimum $L^2$-norm among closed
representatives of the class.  We use the usual Lefschetz decomposition,
Hodge--Riemann bilinear relations, and commutation of the K\"ahler
Lefschetz operator with the Laplacian; see, for example,
\citet{Voisin2002}.

The next estimate concerns a K\"ahler metric stretched in at most half of
the complex tangent directions.  Without its cohomological hypothesis,
a middle-degree norm can grow like $R^{n/2}$.  The cohomological hypothesis forces a uniform loss
in this growth rate. The proof combines pointwise variation of the norm
with orthogonality of harmonic and exact top-degree forms.

\begin{lemma}[A quantitative Hodge norm gap]\label{lem:hodge-gap}
Let $M$ be a compact K\"ahler manifold of complex dimension $2n$, where
$n\ge1$, let $D$ be a K\"ahler form, and let $H$ be a smooth closed
semipositive $(1,1)$-form of pointwise rank at most $n$.  If
$z\in H^{n,n}(M)$ satisfies $[H]^n z=0$, then, for every $R\ge1$,
\begin{equation}\label{hodge:gap}
 \norm{z}_{D+RH}
 \le R^{(n-\delta_n)/2}\norm{z}_{D+H},
 \qquad
 \delta_n=\frac{1}{\binom{2n}{n}+2}.
\end{equation}
\end{lemma}

\begin{proof}
Fix $R$ and write $\omega=D+RH$, $Q=RH$.  Let $\gamma$ be the
$\omega$-harmonic representative of $z$.  In an $\omega$-unitary coframe
diagonalizing $Q$, write its eigenvalues as $b_i\in[0,1]$ and set
$b_I=\sum_{i\in I}b_i$ for an $n$-element subset $I$ of
$\{1,\ldots,2n\}$.  Keeping $\gamma$ fixed while differentiating its
squared norm with respect to $\log R$ gives, relative to the present
volume form,
\begin{equation}\label{hodge:variation}
 \sum_{|I|=|J|=n}
 \bigl(\Tr Q-b_I-b_J\bigr)|\gamma_{I\bar J}|^2.
\end{equation}
Indeed the volume contributes $\Tr Q$, while the covariant holomorphic
and antiholomorphic factors contribute $-b_I$ and $-b_J$.
Define the nonnegative pointwise deficit
\[
 e=n|\gamma|^2-
 \sum_{I,J}(\Tr Q-b_I-b_J)|\gamma_{I\bar J}|^2.
\]
We will prove
\begin{equation}\label{hodge:integrated-deficit}
 \int_M e\,\dd V_\omega\ge\delta_n\norm{\gamma}_\omega^2.
\end{equation}

Put $N_n=\binom{2n}{n}$.  If $\rank Q<n$, then $\Tr Q\le n-1$
and $e\ge|\gamma|^2$.  If $\rank Q=n$, denote its $n$-dimensional
kernel by $K$, the coefficient $\gamma_{K\bar K}$ by $a$, and put
\[
 d_0=n-\Tr Q,\qquad p=\prod_{b_i>0}b_i.
\]
Every coefficient other than $a$ involves a positive eigenvalue in at
least one index set.  Subtracting that eigenvalue from $\Tr Q$ leaves a
sum of at most $n-1$ numbers at most one.  Consequently
\begin{equation}\label{hodge:tensor-deficit}
 e\ge d_0|a|^2+\sum_{(I,J)\ne(K,K)}|\gamma_{I\bar J}|^2.
\end{equation}
Since $0\le b_i\le1$, we also have
\[
 (1-p)^2\le d_0,\qquad 1-p^2\le2d_0.
\]
Consider the top-degree forms
\[
 A=\frac{\omega^n}{n!}\wedge\gamma,
 \qquad B=\frac{Q^n}{n!}\wedge\gamma.
\]
In normalized exterior bases, the coefficient of $A$ is a signed sum of
the $N_n$ diagonal coefficients of $\gamma$.  The coefficient of $B$ is
$p a$, with the same sign as the $a$ term in $A$.  Thus
Cauchy--Schwarz and \eqref{hodge:tensor-deficit} give
\begin{equation}\label{hodge:top-form-bounds}
 |A-B|^2\le N_n e,
 \qquad |B|^2\ge |\gamma|^2-2e.
\end{equation}
The same inequalities hold when $\rank Q<n$: then $B=0$ and
$e\ge|\gamma|^2$.

The form $A$ is harmonic, because multiplication by $\omega$ commutes
with the Laplacian.  The form $B$ is exact, because $Q$ is closed and
\[
 [B]=\frac{R^n}{n!}[H]^n z=0.
\]
They are therefore $L^2$-orthogonal.  Integrating
\eqref{hodge:top-form-bounds} yields
\[
 N_n\int_M e\,\dd V_\omega
 \ge\norm{A-B}_\omega^2
 \ge\norm{B}_\omega^2
 \ge\norm{\gamma}_\omega^2-2\int_M e\,\dd V_\omega,
\]
which is \eqref{hodge:integrated-deficit}.

For completeness, varying the metric does not require differentiating
the harmonic representative.  At the current value of $R$, keeping it
fixed supplies an upper derivative for the minimum defining
$\norm{z}_{D+RH}^2$.  These minimum norms are locally Lipschitz in
$\log R$, by comparison of nearby metrics.  Equations
\eqref{hodge:variation} and \eqref{hodge:integrated-deficit} therefore
imply almost everywhere
\[
 \frac{\dd}{\dd\log R}\norm{z}_{D+RH}^2
 \le(n-\delta_n)\norm{z}_{D+RH}^2.
\]
Integration from $1$ to $R$ proves \eqref{hodge:gap}.
\end{proof}
